
\documentclass[10pt,a4paper]{article}
\usepackage[margin=18mm,headheight=14pt,headsep=8pt,footskip=13mm]{geometry}
\usepackage[T1]{fontenc}
\usepackage{lmodern}
\usepackage{amsmath,amssymb,mathtools}
\usepackage{microtype}
\usepackage[hypertexnames=false]{hyperref}
\usepackage{fancyhdr}
\usepackage{enumitem}
\usepackage{xcolor}
\hypersetup{colorlinks=true,linkcolor=blue!50!black,urlcolor=blue!50!black,pdftitle={Mathematical Appendix: Bayes-Adaptive Market Making},pdfauthor={Hugo Lemonnier}}
\setlength{\parindent}{0pt}
\setlength{\parskip}{4.5pt}
\setlist{nosep,leftmargin=17pt,topsep=3pt}
\emergencystretch=2em
\pagestyle{fancy}\fancyhf{}
\fancyhead[L]{\small Mathematical appendix}
\fancyhead[R]{\small Bayes-adaptive extension}
\fancyfoot[L]{\small Hugo Lemonnier}
\fancyfoot[R]{\small \thepage}
\begin{document}
{\LARGE\bfseries Mathematical appendix\par}
\vspace{2pt}
{\large Joint model-and-regime learning in finite-horizon market making\par}
\vspace{4pt}
Hugo Lemonnier \hfill 27 September 2026\par
\textit{Standalone appendix to the separate Bayes-adaptive research extension.}
\par
\begin{abstract}
This document gives the joint filter, exact finite-horizon recursion, model-only revelation hierarchy, information identities, matched planning interventions, and conservative interpolation and quadrature error allowances for the prespecified two-model synthetic market. Analytic inequalities concern exact expectations. Reported numerical quadrature distances and computed control tables require separate error accounting; no final economic outcome or numerical convergence claim is made here.
\end{abstract}
\vspace{3pt}

This is a separate extension to the core synthetic study. All results below
refer to the specified synthetic model, the finite prior support

\begin{equation*}
M\in\{1,2\},\qquad \theta_1=\theta_2=0.35,\qquad
(\kappa_1,\kappa_2)=(0.002,0.10),\qquad T=30,\quad q_{\max}=2.
\end{equation*}

The model label is fixed within an environment, including its independent pilot
episodes. The initial regime of each episode is drawn independently with equal
sign probabilities. The cold-start model probabilities are equal. Every other
market, execution, accounting, and inventory convention is inherited unchanged
from the task statement. None of the propositions below establishes performance
in a misspecified or real market.

The mathematical quantities in this document are exact expectations and optima.
Implemented tables, computed action scores, and Monte Carlo policy returns are
different objects. Section 8 makes their numerical relationship explicit.

\section{Decision state and chronology}

Let $n=T-t$ be the number of decisions remaining. Immediately before decision
$t$, write

\begin{equation*}
p_{jh}=\Pr(M=j,H_t=h\mid\mathcal I_t),\qquad
w_j=\sum_h p_{jh},\qquad
b_j=\Pr(H_t=+1\mid M=j,\mathcal I_t).
\end{equation*}

The four joint masses are ordered as
$(1,-1),(1,+1),(2,-1),(2,+1)$. When $w_j>0$,
$p_{j,+}=w_jb_j$ and $p_{j,-}=w_j(1-b_j)$. At zero model weight,
the corresponding conditional probability is immaterial. Calculations at that
boundary must not depend on its arbitrary representation.

The sufficient decision state is $(n,q,x,p)$. The following order is essential:

\begin{enumerate}
\item Observe the current public state and select one admissible action.
\item Generate $Z_t$, the return, and selected passive outcomes using \textbf{the
   current} $H_t$. Book actual executions.
\item Observe the return and the outcomes of the submitted quotes. Condition the
   joint belief on those observations to obtain a posterior on $(M,H_t)$.
\item Apply each model's regime transition to obtain the pre-decision belief on
   $(M,H_{t+1})$. Transition the public signal independently and disclose it.
\end{enumerate}

Thus the observation update precedes the hidden-state prediction. A textbook
POMDP convention that observes an emission from the \emph{next} state cannot be
copied into this simulator without changing the indexing.

Cash and absolute midpoint need not appear in the optimization state. If
$W_t=C_t+Q_tP_t$, marked wealth satisfies the accounting identity derived in
Section 3. Current wealth is an additive constant in the remaining objective;
price levels and cash do not affect admissibility, observation laws, or future
increments. The ledger is still required for independent execution verification.
Observed cash and inventory are deterministic functions of the already observed
actions and selected fills and hence supply no additional latent-state evidence.

\section{Exact action-dependent joint filter}

Set $\mu=0.03$, $\sigma=0.30$, and recover the observed innovation as
$z=(r-\mu x)/\sigma$. For a submitted side $s\in\{+1,-1\}$ at depth
$k_s\in\{0,1\}$, define

\begin{equation*}
u_s(x)=\{1+\exp(0.3+0.7k_s+0.2sx)\}^{-1},\qquad
\tau_s(x)=\Phi^{-1}(u_s(x)),
\end{equation*}

\begin{equation}
\ell_h(z,f;x,a)=
\prod_{s\in S(a)}
\Phi\!\left(
 (2f_s-1)\frac{\tau_s(x)-h\theta s z}{\sqrt{1-\theta^2}}
\right).
\tag{1}
\end{equation}

Here $S(a)$ is the set of submitted passive sides and $f$ ranges over their
binary outcomes. The product is one for an empty set. A submitted non-fill is
$f_s=0$; an unsubmitted side contributes no factor. The selected sides are
independent conditional on $(z,h)$, so their factors are multiplied \textbf{before}
mixing over $h$. They are generally not independent after that mixture.

Both models have the same $\theta$, so (1) depends on the model only through
its current regime distribution. The full observation density with respect to
$dz$ and counting measure on $f$ is

\begin{equation}
\phi(z)m(p,z,f),\qquad
m_j=(1-b_j)\ell_-+b_j\ell_+,\qquad
m=\sum_j w_jm_j=\sum_{j,h}p_{jh}\ell_h.
\tag{2}
\end{equation}

The common factor $\phi(z)$ cancels in filtering, but must remain in the
control expectation. The public signal transition similarly cancels from
posterior odds because its known matrix $K$ is model independent.

The posterior on the just-observed current state and the next decision belief
are, respectively,

\begin{equation}
\widetilde p_{jh}=\frac{p_{jh}\ell_h}{m},\qquad
p'_{jh'}=\sum_h\widetilde p_{jh}T_j(h,h'),\qquad
T_j=\begin{pmatrix}1-\kappa_j&\kappa_j\\
\kappa_j&1-\kappa_j\end{pmatrix}.
\tag{3}
\end{equation}

Equivalently,

\begin{equation}
w'_j=\frac{w_jm_j}{m},\qquad
\widetilde b_j=\frac{b_j\ell_+}{m_j},\qquad
b'_j=\kappa_j+(1-2\kappa_j)\widetilde b_j.
\tag{4}
\end{equation}

Denote (3) by $\Psi(p,z,f;x,a)$. For finite $z$, the present model has
strictly positive emission probabilities, so the normalization is positive.
Implementations should nevertheless use log probabilities for filtering and a
defined harmless representation for underflowed zero-mass branches. An exactly
excluded model stays excluded.

For abstention or a market order, $\ell_- =\ell_+=1$. Model weights are
unchanged, and only each conditional regime prediction moves. Public returns
alone contain no evidence about either latent quantity in this model.

\section{Reward and exact Bellman recursion}

Write $h_0=0.025$ for the half spread, $\Delta=0.025$,
$c_P=0.001$, $c_T=0.002$, and $\lambda=0.002$, to avoid confusing
the half spread with the hidden sign $h$. Define
$d_k=h_0+k\Delta$ and $\ell_T(q)=-(h_0+c_T)|q|$.

For a passive action let

\begin{equation*}
\delta q(a,f)=\sum_{s\in S(a)}sf_s,\qquad
c(a,f)=\sum_{s\in S(a)}f_s(d_{k_s}-c_P).
\end{equation*}

For a market order of sign $s$, set $\delta q=s$ and
$c=-(h_0+c_T)$; for abstention both are zero. In all cases,

\begin{equation}
W_{t+1}-W_t=c(a,f)+(q+\delta q)(\mu x+\sigma z).
\tag{5}
\end{equation}

The realized one-period objective increment is (5) minus $\lambda q^2$.
Its conditional expectation is denoted $r(q,x,p,a)$. Let
$\overline h(p)=\sum_{j,h}hp_{jh}=\sum_jw_j(2b_j-1)$. Then

\begin{equation}
\begin{split}
r(q,x,p,a)
={}&q\mu x-\lambda q^2\\
&+\sum_{s\in S(a)}
\left[u_s(x)(d_{k_s}-c_P+s\mu x)
-\sigma\theta\overline h(p)\phi(\tau_s(x))\right]
\end{split}
\tag{6}
\end{equation}

for passive actions, and

\begin{equation}
r(q,x,p,M_s)=(q+s)\mu x-\lambda q^2-h_0-c_T.
\tag{7}
\end{equation}

To verify the covariance term in (6), conditional on $H_t=h$, the pair
$(Z,V_s)$ is standard bivariate normal with correlation $h\theta s$.
Therefore
$E[ZF_s\mid h,x,a]=h\theta s\int_{-\infty}^{\tau_s}v\phi(v)dv
=-h\theta s\phi(\tau_s)$. Multiplication by the inventory sign $s$
gives the last term of (6). The pre-decision inventory penalty is unchanged,
including when both quotes fill.

For any continuation function $F$, define the common-action backup

\begin{equation}
\begin{split}
(\mathcal B F)(q,x,p)
=\max_{a\in\mathcal A(q)}\bigg\{r(q,x,p,a)
 +\sum_f\int_{\mathbb R}\phi(z)m(p,z,f)\,
 \sum_{x'}K_{xx'}
 F(q+\delta q(a,f),x',\Psi(p,z,f;x,a))\,dz\bigg\}.
\end{split}
\tag{8}
\end{equation}

The action set is safe for every possible selected-fill subset. In particular,
an otherwise inadmissible bid at $q=2$ cannot be made safe by also quoting
an ask. Each branch has one future action policy based on its observed joint
posterior. No model-specific maximization occurs inside the model mixture.

\textbf{Proposition 1 (Bayes-adaptive optimum).} The exact optimal value satisfies

\begin{equation}
V^{\mathrm{BA}}_0(q,x,p)=\ell_T(q),\qquad
V^{\mathrm{BA}}_n=\mathcal B V^{\mathrm{BA}}_{n-1}.
\tag{9}
\end{equation}

\emph{Proof.} Equations (1)--(4), inventory accounting, and $K$ give the conditional
law of the next sufficient state from the present sufficient state and action.
Equation (5) telescopes the wealth ledger; the remaining terminal adjustment is
exactly $\ell_T$. Conditional expectation followed by optimization over the
finite admissible action set gives (8). Backward induction proves (9). A
maximizing measurable action exists at every stage; a deterministic tie rule
selects one. Integrability follows from bounded inventory and finite Gaussian
first moments. $\square$

For a fixed admissible policy $\pi$, its actual value $V^\pi_n$ follows the
same recursion with its chosen action in place of the maximum. Consequently,

\begin{equation}
V^\pi_n(q,x,p)\le V^{\mathrm{BA}}_n(q,x,p).
\tag{10}
\end{equation}

This is a Bayesian expectation under the displayed prior $p$, not a claim
that the Bayes-optimal policy dominates every comparator separately in each
true-model environment.

\section{Exact interpretation of posterior-weighted known-model Q values}

Let $V^j_n(q,x,b)$ be the exact optimum when model $j$ is known but the
current regime is hidden. Its filter still uses $\widetilde b$ followed by
the transition with $\kappa_j$. Let $Q^j_n(q,x,b,a)$ force the first
action to be $a$, with optimal \textbf{known-model, hidden-regime} control after
that action. No future shock or future regime is observed by this reference.

Consider an information relaxation that reveals $M$ for free after the
current action and its normal feedback, before the next decision. The regime
remains hidden, including at the revelation time. Revealing the static model
label before or after the independent hidden transition has the same effect
provided the next conditional belief is (4).

\textbf{Proposition 2 (one-action model revelation).} Its optimal objective is

\begin{equation}
U^1_n(q,x,p)=\max_{a\in\mathcal A(q)}
\sum_j w_j Q^j_n(q,x,b_j,a).
\tag{11}
\end{equation}

\emph{Proof.} Before revelation the first action is common to both models. After
normal feedback, the model has posterior probability $w'_j$; if it is
revealed to be $j$, its optimal remaining value is
$V^j_{n-1}(q+\delta q,x',b'_j)$. The branchwise identity
$m w'_j=w_jm_j$ gives

\begin{equation}
\begin{split}
&r(q,x,p,a)+E\!\left[\sum_jw'_j
 V^j_{n-1}(q+\delta q,X',b'_j)\right]\\
&\hspace{20mm}=\sum_jw_jQ^j_n(q,x,b_j,a).
\end{split}
\tag{12}
\end{equation}

Maximizing the common first action proves (11). $\square$

The information available to the relaxed controller contains that of the real
controller. It can ignore its additional label and imitate any admissible
real policy. Combining this observation with (10) gives

\begin{equation}
\boxed{\quad V^\pi_n\le V^{\mathrm{BA}}_n\le U^1_n.\quad}
\tag{13}
\end{equation}

The executable weighted-Q policy selects the maximizing action in (11), updates
its real posterior, and repeats at the next decision. \textbf{Its realized value is
$V^{\pi_{\mathrm{WQ}}}$, not $U^1$}: the promised model revelation never
actually arrives. Weighted-Q can learn the model through its online filter and
can value regime feedback through each $Q^j$. Its planning simplification
is that future \emph{model uncertainty} is removed after the current action.

Equation (13) requires exact optimal known-model Q functions. If those functions
are heuristics or uncertified numerical approximations, the raw computed
maximum is not automatically an upper bound. A regime-observing oracle also
defines a different, more informative comparator than (11).

\section{A hierarchy of delayed-revelation bounds}

For an integer $d\ge0$, reveal the model after exactly $d$ further
decisions, before any later decision. Let $U^d_n$ be the optimum in that
relaxed problem. Revelation at $d=0$ precedes the current decision, so

\begin{equation}
U^0_n(q,x,p)=\sum_jw_jV^j_n(q,x,b_j),\qquad
U^d_0(q,x,p)=\ell_T(q),
\tag{14}
\end{equation}

and for $d\ge1,n\ge1$,

\begin{equation}
U^d_n=\mathcal B U^{d-1}_{n-1}.
\tag{15}
\end{equation}

The delay decreases after each unrevealed decision. Recursing on $U^d_{n-1}$
with the same delay would define a different problem.

\textbf{Proposition 3 (nested information sets).} For all states and $n$,

\begin{equation}
U^0_n\ge U^1_n\ge U^2_n\ge\cdots\ge U^n_n
=U^{n+1}_n=\cdots=V^{\mathrm{BA}}_n.
\tag{16}
\end{equation}

\emph{Proof.} A controller told the model after $d$ actions can ignore it until
after action $d+1$ and implement any policy from the later-revelation
problem. The physical dynamics and permissible actions are the same, proving
each inequality. If $d\ge n$, no remaining decision can use the label;
terminal liquidation is fixed and model independent. Such a label changes no
payoff, proving equality to (9). Equivalently, induction in (15) reaches the
common terminal condition before any useful revelation. $\square$

This proves a hierarchy of \emph{objectives}. It does not prove a monotone ranking
of the real returns of the policies greedy with respect to those objectives.
Their remaining optimism and resulting action errors can differ. Delays one,
two, and three at horizon 30 are limited-lookahead information relaxations;
they are not full Bayes-adaptive optima at horizon 30. At horizon two, delay two
must agree with the exact Bayes-adaptive recursion; at horizon three, delay
three must agree.

Useful exact limiting cases are a singleton model posterior; two identical
models; and $\theta=0$, where the hidden process has no execution or reward
effect. In each case the economically relevant model uncertainty disappears.
For a singleton posterior every delay agrees with the corresponding
known-model hidden-regime optimum.

For a genuine feasible-policy lower value $L\le V^{\mathrm{BA}}$ and a
certified upper value $U\ge V^{\mathrm{BA}}$, $U-L$ bounds the remaining
optimality gap. A simulation estimate of $V^\pi$ and an uncertified
interpolated $U^d$ do not, by themselves, supply such a certificate. Also,
(16) concerns the prior-weighted population. A realized episode or a
true-model-specific mean can exceed a prior-averaged relaxation value without
contradicting the theorem.

\section{What an observation teaches about the model and the regime}

\subsection{The two information quantities share an observation channel}

Condition throughout on the pre-decision information and selected action. With
common $\theta$, the emission depends on $H_t$, not directly on $M$:

\begin{equation}
M\longrightarrow H_t\longrightarrow O_t.
\tag{17}
\end{equation}

This conditional Markov chain implies

\begin{equation}
I(M,H_t;O_t)=I(H_t;O_t)
=I(M;O_t)+I(H_t;O_t\mid M).
\tag{18}
\end{equation}

The two summands give a valid information decomposition, not an additive
decomposition of economic value. Their numerical forms include

\begin{equation}
\mathcal I_M(p,a)
=\sum_f\int\phi(z)m
 \sum_jw'_j\log\frac{w'_j}{w_j}\,dz,
\tag{19}
\end{equation}

\begin{equation}
\mathcal I_{H\mid M}(p,a)
=\sum_jw_j\sum_{h\in\{-1,+1\}}b_{jh}
 \sum_f\int\phi(z)\ell_h\log\frac{\ell_h}{m_j}\,dz,
\quad(b_{j,-},b_{j,+})=(1-b_j,b_j).
\tag{20}
\end{equation}

Zero-weight terms are interpreted by continuity. These are nonnegative
expected information gains. A particular realized posterior entropy can
increase. Entropy computed after predicting $H_{t+1}$ additionally includes
uncertainty introduced by switching and is not the information acquired about
$H_t$. Conditional regime information can be lost through that transition;
the model posterior is unchanged by prediction because the model is static.

The sign-free marginal fill law alone cannot identify either $H_t$ or
$\kappa$ here. This also holds for the two selected side outcomes observed
without their return: simultaneous sign reversal is absorbed by the symmetric
Gaussian innovation, and their joint law depends on $\theta^2$, not the
regime sign. Returns alone are likewise uninformative. Their \textbf{joint}
observation is informative. These statements condition on the known action
and public signal, rather than treating a policy's selected actions as an
independent evidence source.

\subsection{Cold start has exactly zero first-observation model information}

At cold start $b_1=b_2=1/2$. For every action and every observation,
$m_1=m_2$, and hence $w'_j=w_j$ and $\mathcal I_M(p,a)=0$.
More generally this holds whenever the two conditional current-regime beliefs
coincide. The first observation may nevertheless update their common
current-regime posterior $\widetilde b$. Subsequent prediction gives

\begin{equation}
b'_1=0.002+0.996\widetilde b,\qquad
b'_2=0.10+0.80\widetilde b.
\tag{21}
\end{equation}

Unless $\widetilde b=1/2$, those probabilities differ. A later informative
observation can then distinguish persistence models. Initial information
acquisition can prepare a later model-learning opportunity even though its
own immediate model information gain is zero. A one-period pilot episode
therefore supplies no information about $\kappa$ in this two-model family.

The absence of first-step model information is not a reason for the optimal
initial action to match weighted-Q. The policies make different assumptions
about information and uncertainty at future decisions.

\subsection{Complete regime evidence generally cannot be kept while model evidence is erased}

For positive likelihoods, let $R=\ell_+/\ell_-$. The conditional regime
posterior odds multiply by $R$, while the model Bayes factor is

\begin{equation}
\frac{m_1}{m_2}
=\frac{1-b_1+b_1R}{1-b_2+b_2R}.
\tag{22}
\end{equation}

When $b_1\ne b_2$, the right side is strictly monotone in $R$: its
derivative is $(b_1-b_2)/(1-b_2+b_2R)^2$. Thus preserving the full
likelihood-ratio evidence needed by both regime filters also preserves the
model Bayes factor. Except at special states or uninformative observations,
there is no coherent observation deletion that removes model evidence while
retaining all that regime evidence. This is a structural limitation of the
proposed attribution experiment, not a numerical defect.

\section{Matched suppression of model-weight updates in planning}

A useful controlled computational intervention is still available. Starting
from the same actual belief $(w,b_1,b_2)$, compute the true predictive outcome
probability and the true conditional regime updates in (4). Replace only the
next model weights by the current ones when reading a continuation value:

\begin{equation}
\Psi_{\mathrm{FM}}(p,z,f)
=\big(w_j(1-b'_j),\ w_jb'_j\big)_{j=1,2}.
\tag{23}
\end{equation}

For a fixed common continuation $F$, define

\begin{equation}
Q_{F}^{\mathrm{BA}}(a)=r+E_p[F(q',X',\Psi)],\qquad
Q_{F}^{\mathrm{FM}}(a)=r+E_p[F(q',X',\Psi_{\mathrm{FM}})].
\tag{24}
\end{equation}

The two expectations use identical true current outcome probabilities,
immediate rewards, inventory transitions, legal actions, and continuation
function. Only the model-weight update at the next planning boundary changes.
Using $F=V^{\mathrm{BA}}_{n-1}$ makes (24) a direct action-level diagnostic
without changing the downstream continuation function at the same time.

For the executable recursive ablation, let $W_0=\ell_T$ and apply (23) at
every hypothetical step to construct $W_n$. Its policy reads these scores
from the \textbf{true current posterior}, and its online filter uses (3), just as
the other feasible policies do. Thus it suppresses future model-weight updates
in planning, not actual permitted feedback at deployment. Distinguish this
recursive policy comparison from the single-backup experiment in (24).

Important limits are part of its definition:

\begin{itemize}
\item $\Psi_{\mathrm{FM}}$ is an artificial belief intervention, not the Bayes
  posterior for the physical observation experiment. It defines a planning
  surrogate; no upper- or lower-bound claim attaches to $W_n$ or to the
  score difference in (24).
\item The conditional regime filters $b'_j$ are retained, but the marginal
  regime probability changes from $\sum_jw'_jb'_j$ to
  $\sum_jw_jb'_j$. Subsequent forecasts can therefore change along with
  model learning. By (22), a completely clean separation generally cannot be
  realized while preserving all regime evidence.
\item Neither each realized difference in (24) nor its expectation must be
  nonnegative. Jensen's theorem for coherent posterior conditioning does not
  apply to (23).
\item At a state with $b_1=b_2$, (23) agrees with the true update for every
  current observation. The one-backup diagnostic must then be exactly zero
  for any fixed $F$. Recursive BA and frozen-model policies can still
  disagree there because their future continuation functions differ.
\end{itemize}

For comparison, a prediction-only intervention freezes both the observation
update of $w$ and that of each $b_j$, reading continuation at
$b'_j=\kappa_j+(1-2\kappa_j)b_j$. It still integrates the actual stochastic
inventory outcomes. It changes more than the model-weight intervention and
cannot isolate its contribution.

An action-level finding should retain the complete state, admissibility mask,
all action scores under (24), immediate rewards, numerical refinements, and
model/regime information diagnostics. An action difference stable at relevant
numerical resolutions supports sensitivity to the declared planning update.
Fresh paired objective comparisons are needed to establish economic value.
Lower entropy, a selected action disagreement, or a negative immediate reward
alone establishes neither model-learning benefit nor profitable exploration.
If the specified contrast is negligible, adverse, or numerically unstable,
that is the finding; it is not a reason to relabel regime information as model
information or to select a new preferred hypothesis after evaluation.

\section{Numerical error: what can be proved and what remains empirical}

\subsection{Convexity and the product coordinates}

For a fixed $(n,q,x)$ and a fixed nonanticipative policy tree, the expected
remaining payoff is linear in the initial four-state mass vector $p$.
Taking the supremum over the same class of admissible policy trees proves
convexity of $V_n^{\mathrm{BA}}$ in $p$. The same argument holds for
$U^d_n$, where the policy tree also has its prescribed future revelation.
Continuous observations mean that this need not be a maximum of finitely many
linear functions; a finite latent state space alone does not justify a finite
piecewise-linear representation.

On a rectangular grid in $(w,b_1,b_2)$, with $w=w_1$, the mapping

\begin{equation}
p(w,b_1,b_2)=\big(w(1-b_1),wb_1,(1-w)(1-b_2),(1-w)b_2\big)
\tag{25}
\end{equation}

is multiaffine. The nonnegative tensor interpolation weights $\alpha_v$ on
the eight corners satisfy
$p(w,b_1,b_2)=\sum_v\alpha_vp_v$. Therefore, for exact nodal values of
any convex joint-belief value $V$,

\begin{equation}
V(p)\le\sum_v\alpha_vV(p_v).
\tag{26}
\end{equation}

This fact concerns joint-mass barycenters; convexity in the three conditional
coordinates is not assumed. With exact integration, exact maximization, and
nonnegative interpolation, an induction from the terminal payoff makes the
gridded BA calculation an upper approximation. Positive quadrature weights
preserve monotonicity but \textbf{do not give a sign to Gaussian quadrature error}.
Thus (26) alone does not certify the implemented Gauss--Hermite tables.

\subsection{A global payoff and interpolation envelope}

A conservative exact bound for any conditional expected stage reward is
$|r(q,x,\delta_{jh},a)|\le r_*=0.231$. To check it, $|q\mu x|\le0.06$,
$\lambda q^2\le0.008$, $u_s<1/2$,
$|d_k-c_P+s\mu x|\le0.079$, and
$\sigma\theta\phi(\tau_s)<0.042$. There are at most two submitted sides.
Market actions have a smaller bound. Terminal liquidation lies in
$[-0.054,0]$. Consequently every policy tree's expected payoff, conditional
on any initial hidden state, lies in

\begin{equation*}
[-0.231n-0.054,\ 0.231n].
\end{equation*}

Let $D_n=0.231n+0.027$, half this common payoff envelope's width. Subtracting
its midpoint from each linear policy coefficient proves

\begin{equation}
|V_n(p)-V_n(\bar p)|\le D_n\|p-\bar p\|_1
\tag{27}
\end{equation}

for exact BA and fixed-delay revelation values at fixed $(q,x)$.
This uses a bound on \textbf{conditional expected} stage rewards, not an incorrect
bounded-support assertion about Gaussian realized PnL.

If the containing cell has local widths $\Delta_w,\Delta_{b_1},\Delta_{b_2}$, the tensor corner distribution gives

\begin{equation}
\sum_v\alpha_v\|p_v-p\|_1
\le \Delta_w+w\Delta_{b_1}+(1-w)\Delta_{b_2}
\le\Delta_w+\max_j\Delta_{b_j}.
\tag{28}
\end{equation}

For example, split the difference by first changing $w$ and then $b_1,b_2$.
The respective $\ell_1$ changes are $2|\delta w|$ and
$2[w|\delta b_1|+(1-w)|\delta b_2|]$. Linear endpoint sampling has expected
absolute displacement at most half its interval width, proving (28).
Equations (26)--(28) bound interpolation error by
$D_n(\Delta_w+\max_j\Delta_{b_j})$.

\subsection{A global Gaussian quadrature envelope without dividing by a tiny posterior mass}

Set

\begin{equation*}
c_\theta=\frac{\theta}{\sqrt{1-\theta^2}\sqrt{2\pi}}<0.15.
\end{equation*}

Each conditional Bernoulli fill probability is $c_\theta$-Lipschitz in
$z$. For $s_a\le2$ submitted sides, the $\ell_1$ distance between
the two product Bernoulli outcome distributions at $z$ and $\bar z$
is at most $2s_ac_\theta|z-\bar z|$.

For each fill branch, use its \textbf{unnormalized} next-state vector
$v_f(z)_{jh'}=\sum_hp_{jh}\ell_h(z,f)T_j(h,h')$. It has mass $m$.
The homogeneous value perspective
$mV(q',x',v_f/m)$ is the supremum of linear policy payoffs in $v_f$.
After subtracting the common payoff-envelope midpoint times the mass, it is
$D_n$-Lipschitz in $v_f$. Hidden-state prediction contracts
$\ell_1$ distance. Summing over fills and averaging over $K$ therefore
shows that the exact continuation integrand

\begin{equation*}
g(z)=\sum_fm(p,z,f)\sum_{x'}K_{xx'}V_n(q',x',\Psi(p,z,f))
\end{equation*}

is $2s_ac_\theta D_n$-Lipschitz. The midpoint terms cancel because
$\sum_fm(p,z,f)=1$ at every $z$. This proof avoids any lower bound on
individual posterior denominators.

Let $\nu=\sum_i\omega_i\delta_{z_i}$ be a normalized positive numerical
quadrature law and let
$\mathcal W_1=W_1(\mathcal N(0,1),\nu)$. Coupling the two scalar variables
gives the deterministic inequality

\begin{equation}
\left|E[g(Z)]-\sum_i\omega_i g(z_i)\right|
\le 2s_ac_\theta D_n\mathcal W_1
\le 4c_\theta D_n\mathcal W_1.
\tag{29}
\end{equation}

This applies even though continuation has decision kinks. It requires no
unjustified high-order derivative bound for Gauss--Hermite quadrature.

The transport distance has a one-dimensional exact expression. Sort the nodes,
put $c_i=\sum_{k\le i}\omega_k$, and
$a_i=\Phi^{-1}(c_i)$, with $a_0=-\infty$ and $a_N=+\infty$. Then

\begin{equation}
\mathcal W_1=\sum_{i=1}^N
\int_{a_{i-1}}^{a_i}|z-z_i|\phi(z)\,dz.
\tag{30}
\end{equation}

Each term can be evaluated by splitting at $z_i$ and using
$\int z\phi(z)dz=-\phi(z)$. Floating evaluation for ideal normalized
Gauss--Hermite rules gives approximately 0.2020035, 0.1567188, 0.1224879, and
0.1004662 for 15, 25, 41, and 61 nodes. These displayed evaluations are not
outward-rounded interval certificates for the stored floating-point rule.

\subsection{Propagation, one-sided control, and the size of the allowance}

Bellman maximization and expectation are nonexpansive in the uniform norm.
With one grid, quadrature rule, exact analytic rewards, and exact arithmetic,
let $\delta_G$ be the maximum over all grid cells of $\Delta_w+\max_j\Delta_{b_j}$. This includes nonuniform meshes. Starting from the exact terminal
value, the error at computed grid nodes obeys

\begin{equation}
\max_{q,x,p\ \mathrm{on\ grid}}
|\widehat V_n(q,x,p)-V_n^{\mathrm{BA}}(q,x,p)|
\le\sum_{k=0}^{n-1}D_k(\delta_G+4c_\theta\mathcal W_1).
\tag{31}
\end{equation}

To see this, interpolate the exact preceding-horizon nodal values, apply (28),
and add (29). The previous uniform nodal error passes through the positive
normalized operator without amplification. The maximum over actions cannot
increase that error. Off-grid value queries can require an additional
$D_n\delta_G$ for the last interpolation; this term must not be omitted
from a claimed two-sided guarantee.

There is a useful one-sided improvement: interpolation of the exact value
never lowers it. At every backup, compare the computed continuation to the
exact value using (26) before applying (29). Induction proves, including for
an interpolated final value,

\begin{equation}
V_n^{\mathrm{BA}}(q,x,p)
\le\widehat V_n(q,x,p)+
\sum_{k=0}^{n-1}4c_\theta D_k\mathcal W_1.
\tag{32}
\end{equation}

The bound has no interpolation penalty on this side. Analogous arguments
apply to delayed-revelation recursions if the terminal known-model values
carry their own applicable error allowance. They must not be assumed exact
merely because they came from the frozen known-model solver.

These formulas are a route to a conservative certificate, not evidence that a
tight certificate has been implemented. Validating floating normalization,
special-function evaluations, all arithmetic, and tie tolerance adds its own
per-stage allowance. A chosen action whose score is within the $10^{-12}$
tie tolerance of the maximum adds at most that tolerance per exact backup;
floating arithmetic requires a separately justified allowance.

The magnitude matters. At $n=30$, $\sum_{k=0}^{29}D_k=101.295$.
Using the displayed 61-node transport distance makes the allowance in (32)
about \textbf{6.07 objective units}, before floating-point certification. For illustration only, a uniform 33-by-65-by-65 grid has $\delta_G=0.046875$, adding about 4.75 units to the two-sided nodal envelope (31). Any other mesh uses its own maximum cell-width sum. These safe-envelope calculations are far too
loose to certify a 0.002 economic gap. Their existence does not justify
calling a much smaller observed refinement difference a rigorous error bound.

For an action choice, a uniform $\varepsilon$ bound on every action score
would certify a unique leading action when its estimated lead exceeds
$2\varepsilon$. It would bound the chosen action's true one-step loss by
$2\varepsilon$, and such losses could be accumulated over a finite horizon.
No such small uniform $\varepsilon$ follows from a few checked states.

The practical validation must therefore report grid and quadrature refinement,
short-horizon exact-recursion checks, singleton-model agreement, reached-state
action stability, and independently integrated Bellman residuals as \textbf{numerical
evidence}. Pointwise residuals do not control unsampled states. Monte Carlo
intervals address sampling variability of an executable policy under the
evaluation population; they do not cover solver bias or supply an optimality
gap. Plots should identify $\widehat U^d$ as numerical relaxation estimates,
separately from achieved-policy means and their sampling intervals.

\section{Pilot episodes, reset semantics, and validation identities}

For independent pilot episodes $D_1,\ldots,D_E$ from the same fixed model,

\begin{equation}
w_j(D_{1:E})\propto w_j^{\mathrm{prior}}
\prod_{e=1}^E L_j(D_e),\qquad
L_j(D_e)=\sum_{h_0,\ldots,h_{T-1}}
\frac12\left[\prod_{t=0}^{T-1}\ell_{h_t}(z_t,f_t)\right]
\left[\prod_{t=0}^{T-2}T_j(h_t,h_{t+1})\right].
\tag{33}
\end{equation}

Common return and signal densities are omitted only because they cancel in
model odds. The pilot action probabilities also cancel: the behavior policy
is known and chooses uniformly among publicly admissible actions. Conditioning
on the realized actions does not license access to unselected fills.

Each episode starts the conditional regime filter at $b_j=1/2$.
Accumulate its model log evidence, retain model weights across pilot episodes,
and reset regime beliefs at the next independent episode. Carrying the final
regime posterior into the next independent pilot fabricates persistence;
resetting model weights discards valid shared-model evidence.

Every fresh evaluation episode starts at
$q=0,x=0,p_{jh}=w_j(D_{1:E})/2$. If the declared estimand is performance
after a fixed pilot budget, evaluation episodes independently reuse that same
pilot posterior. They do not gradually increase each other's training budget.
All feasible policies receive the identical pilot within a replicate. Model
labels and exogenous tapes belong to the evaluator, not the policy interface.

Useful independent checks of the derivation are:

\begin{itemize}
\item Joint prediction equals (3), and its model marginal equals (4).
\item All selected-fill branches integrate to one; market/abstention updates
  change no model weights; the unused conditional belief at zero weight has
  no effect.
\item A one-period evaluation uses $\ell_T(q')$, not a second liquidation;
  post-terminal hidden prediction cannot change its value. Every belief-update
  suppression has identical one-period scores when it shares the actual
  inventory outcomes, because the terminal continuation is independent of
  belief.
\item Singleton-model BA values agree with the known-model hidden-regime
  recursion. Identical models collapse in the same way.
\item Delay one equals weighted exact known-model Q action by action;
  delay $n$ agrees with BA at short horizon $n$.
\item Model information and the single-backup frozen-weight score difference are
  exactly zero when $b_1=b_2$, including cold start.
\item The online policy value, the numerical Bellman table, and a model-revelation
  objective are never substituted for one another in reports.
\end{itemize}

These checks belong to the predeclared E2E validation artifact. They are not a
request to add post-implementation unit tests to the frozen core.

\clearpage
\section{Attribution to established work}

The Bayes-adaptive idea---include uncertainty about the environment model in the
hidden state and plan jointly for model learning, state identification, and
reward---is established. Ross, Chaib-draa, and Pineau formulate it for uncertain
transition and observation probabilities using count-augmented states. The
present two-point prior is a finite, simpler instance of the same decision
principle; it does not require their Dirichlet construction.
\href{https://papers.nips.cc/paper_files/paper/2007/file/3b3dbaf68507998acd6a5a5254ab2d76-Paper.pdf}{Ross, Chaib-draa and Pineau (2007), \emph{Bayes-Adaptive POMDPs}, Sections 2--3}.

The use of extra information to upper-bound partially observed control and
the distinction between an approximate value and the control policy extracted
from it are also established. Hauskrecht analyzes QMDP, bound hierarchies, and
convex interpolation. Our weighted-Q interpretation is related to QMDP but
reveals \textbf{only the model}, while its regime remains hidden. Equations
(11)--(16) explicitly derive the appropriate information relaxation for this
task instead of importing the fully state-observing QMDP formula.
\href{https://jair.org/index.php/jair/article/download/10262/24449}{Hauskrecht (2000), \emph{Value-Function Approximations for Partially Observable Markov Decision Processes}, Sections 3.2, 4.1--4.2, 4.5}.

Smallwood and Sondik establish finite-horizon belief-space control and the
convex piecewise-linear structure in their finite formulation. We use the
policy-as-linear-payoff argument to establish convexity directly for this
task's continuous return observation, without claiming finitely many value
pieces. The filter, accounting specialization, model-only revelation proof,
delayed hierarchy, common-emission attribution limitation, and error envelopes
above are worked derivations for this particular extension; no general-method
novelty is claimed.
\href{https://pubsonline.informs.org/doi/10.1287/opre.21.5.1071}{Smallwood and Sondik (1973), \emph{The Optimal Control of Partially Observable Markov Processes over a Finite Horizon}}.
\end{document}
